\section{领域五：上下文管理与可靠性（15\%）}

权重最小的领域，但原文强调"这里的错误会级联到所有其他地方"。

\subsection{渐进式摘要的陷阱}

\begin{warningbox}{渐进式摘要会丢失关键事务数据}
Progressive summarisation（渐进式摘要）会丢失金额、日期、订单号等事务性数据。正确做法是维护一个持久的 \textbf{"case facts" 块}，包含所有提取的关键数据字段，\textbf{永不摘要}，并在每次 prompt 中包含。
\end{warningbox}

\subsection{"Lost in the Middle" 效应}

\begin{importantbox}{关键信息放在开头}
模型容易遗漏埋在长输入中间的内容（"lost in the middle" 效应）。解决方案：将关键摘要和重要发现放在输入的\textbf{开头}位置。
\end{importantbox}

\subsection{升级触发器（Escalation Triggers）}

\begin{table}[H]
\centering
\begin{tabular}{lp{7cm}c}
\toprule
\textbf{触发器} & \textbf{描述} & \textbf{可靠性} \\
\midrule
客户主动请求人工 & 立即执行，不附加条件 & 可靠 \\
策略空白 & Agent 遇到规则未覆盖的情况 & 可靠 \\
无法推进 & Agent 无法继续完成任务 & 可靠 \\
\midrule
情绪分析 & 基于客户情绪判断是否升级 & \textbf{不可靠} \\
自报置信度分数 & 模型自己报告的置信度 & \textbf{不可靠} \\
\bottomrule
\end{tabular}
\caption{升级触发器的可靠性对比}
\end{table}

\begin{warningbox}{考试陷阱：不可靠的升级触发器}
情绪分析（sentiment analysis）和模型自报置信度分数（self-reported confidence scores）是考试中常见的干扰选项。它们看起来合理，但实际上不够可靠，不应作为升级触发条件。
\end{warningbox}

\subsection{错误传播的正确方式}

正确的错误传播应包含结构化上下文：
\begin{itemize}[leftmargin=1.5em]
    \item 失败类型（failure type）
    \item 尝试的查询（attempted query）
    \item 部分结果（partial results）
    \item 替代方案（alternatives）
\end{itemize}

两种必须避免的反模式：
\begin{itemize}[leftmargin=1.5em]
    \item 静默吞掉错误（silently suppressing errors）
    \item 因单个失败而终止整个工作流（killing entire workflows on single failures）
\end{itemize}

\subsection{推荐实践项目}

\begin{knowledgebox}{动手练习：错误传播与部分结果}
构建一个包含两个子 Agent 的 coordinator 系统：
\begin{itemize}[leftmargin=1.5em]
    \item 模拟一个子 Agent 超时
    \item 验证 coordinator 能收到结构化的错误上下文
    \item 确认 coordinator 能基于部分结果继续推进
    \item 测试面对冲突信息源时的处理逻辑
\end{itemize}
\end{knowledgebox}

\subsection{本章小结}

上下文管理的关键要点：用持久化的 "case facts" 块替代渐进式摘要来保留事务数据；利用 "lost in the middle" 效应的规律将关键信息前置；仅使用可靠的升级触发器（客户请求、策略空白、无法推进）；错误传播要结构化，既不静默吞掉也不因单点失败而全面崩溃。

